\newpage
\section*{Apéndices}
\addcontentsline{toc}{section}{Apéndices}

Este capítulo lista los artefactos anexos activos y aclara el estado de aquellos elementos que todavía dependen del freeze final del producto.

\subsection*{Apéndice A: Repositorio de código fuente}
\addcontentsline{toc}{subsection}{Apéndice A: Repositorio de código fuente}

Repositorio principal del proyecto:
\begin{center}
\url{https://github.com/delarge95/WebGL-Thesis-Proposal}
\end{center}

\noindent Este repositorio es la fuente de verdad técnica para escena, scripts, herramientas de editor y documentación de desarrollo.

\subsection*{Apéndice B: Manual de usuario}
\addcontentsline{toc}{subsection}{Apéndice B: Manual de usuario}

Ruta del manual activo:
\begin{center}
\texttt{Informe\_final/Manual\_de\_usuario/manual\_usuario.tex}
\end{center}

\subsection*{Apéndice C: Manual técnico}
\addcontentsline{toc}{subsection}{Apéndice C: Manual técnico}

Ruta del manual activo:
\begin{center}
\texttt{Informe\_final/Manual\_tecnico/manual\_tecnico.tex}
\end{center}

\subsection*{Apéndice D: Documentación técnica de arquitectura}
\addcontentsline{toc}{subsection}{Apéndice D: Documentación técnica de arquitectura}

Documentos activos:
\begin{itemize}
    \item \texttt{Desarrollo\_App/Documentacion\_Tecnica/01\_Arquitectura\_del\_Sistema.md}
    \item \texttt{Desarrollo\_App/Documentacion\_Tecnica/02\_Referencia\_Tecnica\_Modulos.md}
\end{itemize}

\subsection*{Apéndice E: Matriz de desconexiones}
\addcontentsline{toc}{subsection}{Apéndice E: Matriz de desconexiones}

La matriz maestra de desconexiones entre app y documentación se encuentra en:
\begin{center}
\texttt{Desarrollo\_App/MATRIZ\_DESCONEXIONES\_APP\_DOCUMENTACION.md}
\end{center}

\subsection*{Apéndice F: Catálogo canónico del Holybro X500 V2}
\addcontentsline{toc}{subsection}{Apéndice F: Catálogo canónico del Holybro X500 V2}

La taxonomía semántica del proyecto se fija a partir de:
\begin{center}
\texttt{desarrollo/docs/investigacion/Holybro/x500v2\_parts\_data.json}
\end{center}

\noindent Esta fuente define 28 piezas canónicas. La escena final opera adicionalmente con 30 anchors y 257 renderers/colliders auditados.

\subsection*{Apéndice G: Validación funcional de la escena final}
\addcontentsline{toc}{subsection}{Apéndice G: Validación funcional de la escena final}

El documento de referencia para la build funcional es:
\begin{center}
\texttt{Desarrollo\_App/VALIDACION\_FUNCIONAL\_FINA\_2026-04-09.md}
\end{center}

\subsection*{Apéndice H: Instrumentos de evaluación}
\addcontentsline{toc}{subsection}{Apéndice H: Instrumentos de evaluación}

Los instrumentos metodológicos activos del cierre son:
\begin{itemize}
    \item cuestionario SUS;
    \item NASA-TLX Raw;
    \item protocolo \textit{Think-Aloud};
    \item tablas de KPIs técnicos a diligenciar post-freeze.
\end{itemize}

\noindent Las tareas de prueba deben alinearse con el flujo visible final: navegar, seleccionar pieza, revisar el \textit{bottom sheet}, usar \texttt{Inspect}, aplicar \texttt{Explode} o \texttt{Cut}, y cambiar de modo visual en \texttt{Studio}.

\subsection*{Apéndice I: Estado del pipeline de activos}
\addcontentsline{toc}{subsection}{Apéndice I: Estado del pipeline de activos}

Para efectos de cierre documental:
\begin{itemize}
    \item Blender y Unity conforman la ruta defendible de la versión final activa;
    \item Houdini se mantiene como herramienta evaluada o experimental, no como pipeline oficial de la build cerrada, salvo evidencia nueva derivada del reimport final;
    \item la optimización CAD/import pendiente debe completarse antes del freeze definitivo.
\end{itemize}

\subsection*{Apéndice J: Demo público}
\addcontentsline{toc}{subsection}{Apéndice J: Demo público}

La URL pública definitiva del demo debe registrarse solo cuando exista una build congelada, publicada y verificada. Este apéndice no fija una URL transitoria para evitar inconsistencias entre documento y producto.

\subsection*{Apéndice K: Audio y otros pendientes}
\addcontentsline{toc}{subsection}{Apéndice K: Audio y otros pendientes}

Al momento de este cierre documental:
\begin{itemize}
    \item no existen clips de audio integrados en \texttt{Assets};
    \item por defecto, el audio se clasifica como trabajo futuro si no entra antes del freeze;
    \item la evaluación de usabilidad y el capítulo 5 siguen condicionados a la build final optimizada.
\end{itemize}

\subsection*{Apéndice L: Auditoría matemática y trazabilidad con Wolfram}
\addcontentsline{toc}{subsection}{Apéndice L: Auditoría matemática y trazabilidad con Wolfram}

La auditoría activa de coherencia entre fórmulas documentadas, lógica de código y arquitectura real se consolida en:
\begin{center}
\texttt{Desarrollo\_App/AUDITORIA\_MATEMATICA\_Y\_ARQUITECTURA\_2026-04-12.md}
\end{center}

\noindent Como soporte técnico adicional, el registro histórico de verificaciones del subsistema térmico mediante el workflow local basado en WolframAlpha se conserva en:
\begin{center}
\texttt{desarrollo/docs/sistema\_termico/wolfram\_verificaciones.md}
\end{center}

\noindent Ambos documentos respaldan la interpretación correcta del modo \texttt{Thermal} como simulación híbrida heurística y no como FEA en tiempo real.

\newpage
